\documentclass[10pt,twocolumn]{article}

\usepackage[margin=2.2cm]{geometry}
\usepackage{amsmath,amssymb,amsfonts}
\usepackage{booktabs}
\usepackage{graphicx}
\usepackage[colorlinks=true,citecolor=blue,linkcolor=blue,urlcolor=blue]{hyperref}
\usepackage{natbib}
\usepackage{microtype}
\usepackage{xcolor}
\usepackage{caption}
\usepackage{enumitem}
\usepackage{array}

\setlength{\parskip}{0.4em}
\setlist{nosep,leftmargin=1.2em}

\title{Anonymization Under Attack: Evaluating PII Protection\\Beyond Detection Accuracy}
\author{
  Mr A.\thanks{International Institute of Information Technology, Hyderabad.}
  % TODO: add remaining team member names here
  \\[2pt]
  \texttt{huggingface.co/datasets/huggingbahl21/saha-al}
}
\date{}

\begin{document}
\maketitle

% ============================================================
\begin{abstract}
Text anonymization systems are routinely evaluated on detection accuracy or token-level recall, yet these metrics ignore whether anonymized text resists adversarial re-identification.
We introduce SAHA-AL, a benchmark that evaluates PII anonymization as a \emph{system under attack}, measuring privacy risk from retrieval-based, statistical, and generative adversaries alongside a formalized privacy-utility tradeoff.
Across eight baselines, fine-tuned seq2seq models reduce entity leakage from 26--83\% (rule-based) to below 1\%, but retrieval attacks still recover 1.9\% of entities---$\approx$15$\times$ more than generative LLM attacks.
Standard evaluation metrics fail to expose this residual risk: anonymization benchmarks must incorporate adversarial threat models.
\end{abstract}

% ============================================================
\section{Introduction}\label{sec:intro}

Text anonymization is a prerequisite for privacy-preserving data sharing in healthcare, legal, and financial domains.
In practice, organizations deploy rule-based pipelines~\citep{presidio2023,spacy2020} or fine-tuned models~\citep{lewis2020bart,raffel2020t5} to replace personally identifiable information (PII) with synthetic alternatives, then assume the output is safe to release.
This assumption is rarely tested.

Existing benchmarks evaluate anonymization as a standalone text transformation: detection $F_1$, token recall, or surface-level leakage rates.
None systematically subject the anonymized output to adversarial attack.
This creates a dangerous gap: a system may appear to achieve low leakage while remaining vulnerable to embedding-based retrieval~\citep{shokri2017} or LLM-based inference~\citep{staab2024}.
Surface anonymization is not semantic privacy.

SAHA-AL closes this gap.
The name refers to the broader project: a \textbf{S}emi-\textbf{A}utomated \textbf{H}uman-in-the-loop \textbf{A}nnotation pipeline with \textbf{A}ctive \textbf{L}earning that produces the gold-standard dataset, on top of which this benchmark is built.
The benchmark itself decomposes evaluation into three tasks---detection, anonymization quality, and adversarial privacy risk---and defines a three-tier threat model covering passive observation, retrieval-based re-identification, and generative inference.
The core thesis of this work is that \textbf{PII anonymization must be evaluated as a system under adversarial attack, not as a standalone text transformation task.}

\paragraph{Contributions.}
\begin{enumerate}
\item A multi-task evaluation framework with 11~metrics spanning leakage, utility, format preservation, and adversarial attack resistance.
\item A three-tier adversarial threat model incorporating statistical (CRR-3), retrieval (ERA), and generative (LRR) attacks.
\item A parameterized Privacy-Utility Score (PUS) and Pareto frontier analysis formalizing the privacy-utility tradeoff.
\item A diagnostic failure taxonomy that classifies entity-level errors into seven categories (five of which are active on the tested systems).
\item Empirical evidence on eight baselines showing that (a)~end-to-end seq2seq models reduce leakage by an order of magnitude over rule-based pipelines, and (b)~retrieval attacks expose residual privacy risk invisible to standard metrics.
\end{enumerate}

% ============================================================
\section{Related Work}\label{sec:related}

\paragraph{PII detection and anonymization benchmarks.}
Prior PII benchmarks focus predominantly on detection accuracy.
The Text Anonymization Benchmark (TAB)~\citep{pilan2022tab} evaluates entity-level masking on European court decisions but does not measure anonymization quality or adversarial risk.
Clinical de-identification shared tasks~\citep{stubbs2015} evaluate span-level detection on medical notes via $F_1$, treating anonymization as a detection-only problem.
\citet{lison2021} survey text anonymization methods and note that most evaluations measure only precision and recall of entity detection without assessing whether anonymized outputs resist re-identification.
SAHA-AL differs by evaluating anonymization as a \emph{system under attack}: detection is one of three tasks, complemented by leakage, utility, and adversarial resistance metrics.

\paragraph{Adversarial privacy attacks.}
Membership inference attacks~\citep{shokri2017} assess whether a record was in the training data but do not target anonymized text.
\citet{staab2024} demonstrate that LLMs can infer personal attributes from text, motivating our LRR metric.
Embedding-based re-identification has been explored in record linkage~\citep{narayanan2008} but not previously applied as a benchmark metric for text anonymization.
Our ERA metric operationalizes retrieval-based re-identification as a standardized evaluation component.

\paragraph{Anonymization methods.}
Rule-based systems such as Presidio~\citep{presidio2023} and spaCy NER pipelines~\citep{spacy2020} follow a detect-then-replace paradigm.
Neural approaches include fine-tuned sequence-to-sequence models~\citep{lewis2020bart,raffel2020t5} that learn end-to-end text rewriting, and hybrid pipelines combining NER with conditional generation.
Differential privacy~\citep{dwork2006} provides formal guarantees but is typically applied to aggregate statistics rather than free-text anonymization.

% ============================================================
\section{Problem Formulation}\label{sec:formulation}

Let $\mathcal{D} = \{(x_i, E_i)\}_{i=1}^{N}$ be a dataset where $x_i$ is a text document and $E_i = \{e_{i,1}, \ldots, e_{i,k_i}\}$ is the set of PII entities, each $e = (\text{text}, \text{type}, \text{start}, \text{end})$ with surface text, type label, and character-level span offsets.

An \emph{anonymization system} $f$ produces:
\begin{equation}
  \hat{x}_i = f(x_i)
\end{equation}
where $\hat{x}_i$ is the anonymized text with entities replaced by synthetic alternatives.
A \emph{detection system} $g$ produces predicted spans $\hat{E}_i = g(x_i) = \{(\text{start}_j, \text{end}_j, \text{type}_j)\}_{j=1}^{m_i}$.

\paragraph{Threat model.}
We consider three adversary classes, in order of increasing capability:
\begin{enumerate}
\item \textbf{Passive observer.} Has access only to $\hat{x}_i$. Can the original entity be read directly from the output? Measured by ELR.
\item \textbf{Knowledge-enriched adversary.} Has access to a database of candidate entities (e.g., from training data or a census). Can they recover the original entity via embedding similarity? Measured by ERA@$k$.
\item \textbf{Generative adversary.} Has access to a powerful language model. Can they infer the original entity from contextual clues in $\hat{x}_i$? Measured by LRR.
\end{enumerate}

This hierarchy reflects realistic deployment threats: passive leakage is the minimum bar, but organizations must also defend against adversaries with auxiliary knowledge or generative capabilities.

% ============================================================
\section{Benchmark Design}\label{sec:benchmark}

This section describes the core contribution: an evaluation framework that treats anonymization as a system evaluated under multiple adversarial capabilities.

\subsection{Task 1: PII Detection}

\textbf{Input:} Original text $x_i$. \textbf{Output:} Detected entity spans $\hat{E}_i$.
Predicted spans are matched to gold spans via greedy bipartite matching in three modes: \emph{exact} (identical offsets), \emph{partial} (IoU~$>0.5$), and \emph{type-aware} (exact match plus type agreement).
Corpus-level micro-averaged precision, recall, and $F_1$ are reported.

\subsection{Task 2: Text Anonymization Quality}

\textbf{Input:} Original text $x_i$. \textbf{Output:} Anonymized text $\hat{x}_i$.

\paragraph{Entity Leakage Rate (ELR).}
The fraction of gold entities whose text appears verbatim in $\hat{x}_i$, using word-boundary case-insensitive regex:
\begin{equation}
\text{ELR} = \frac{\sum_{i}\sum_{j} \mathbb{1}[e_{i,j} \in \hat{x}_i]}{|\mathcal{E}|} \times 100\%.
\end{equation}

\paragraph{Token Recall.}
Fraction of entity-span tokens (length~$\geq 2$) absent from $\hat{x}_i$. Higher is better.

\paragraph{Over-Masking Rate (OMR).}
Fraction of non-entity tokens unnecessarily altered. Lower is better.

\paragraph{Format Preservation Rate (FPR).}
For structured types (EMAIL, PHONE, SSN, CREDIT\_CARD, ZIPCODE, DATE), the fraction of replacements matching an expected format regex.

\paragraph{BERTScore $F_1$.}
Semantic similarity between $x_i$ and $\hat{x}_i$ via \texttt{distilbert-base-uncased} embeddings~\citep{zhang2020bertscore}, averaged across records.

\subsection{Task 3: Privacy Risk Assessment}

This task subjects anonymized outputs to adversarial attack.
No analogous evaluation exists in prior PII benchmarks.

\paragraph{CRR-3 (Capitalized 3-gram Survival Rate).}
Fraction of capitalized 3-grams from $x_i$ appearing in $\hat{x}_i$. Captures named-entity structure survival beyond individual entity matching.

\paragraph{ERA@$k$ (Entity Recovery Attack).}
A retrieval adversary encodes candidate entities (per type, from training data, capped at 500) and anonymized texts with Sentence-BERT~\citep{reimers2019}.
For each gold entity, candidates are ranked by cosine similarity to the anonymized text embedding.
ERA@$k$ is the fraction recovered in the top-$k$:
\begin{equation}
\text{ERA@}k = \frac{|\{e : e.\text{text} \in \text{Top-}k(\cos(\hat{x}, \text{pool}_{e.\text{type}}))\}|}{|\mathcal{E}_{\text{eval}}|} \times 100\%.
\end{equation}

\paragraph{LRR (LLM Re-identification Rate).}
A generative adversary is prompted to guess original entities from anonymized text.
Exact and fuzzy ($>0.8$ character similarity) match rates are reported.
Evaluated on 300 sampled records using Qwen~2.5-14B and Llama~3.3-70B.

\paragraph{UAC (Unique Attribute Combination Rate).}
Fraction of records with a unique surviving quasi-identifier combination, approximating $k$-anonymity violation at $k\!=\!1$~\citep{sweeney2002}.

\subsection{Privacy-Utility Tradeoff}

We define the Privacy-Utility Score:
\begin{equation}
\text{PUS}(\lambda) = \lambda \cdot \underbrace{\left(1 - \tfrac{\text{ELR}}{100}\right)}_{\text{Privacy}} + (1-\lambda) \cdot \underbrace{\tfrac{\text{BERTScore}}{100}}_{\text{Utility}}
\end{equation}
where $\lambda \in [0,1]$.
A system is \emph{Pareto-optimal} if no other system achieves both strictly higher privacy and strictly higher utility.
PUS provides a single scalar for system comparison at any operating point.

\subsection{Failure Taxonomy}

The implementation classifies each entity replacement into one of seven categories, evaluated in priority order: \emph{Full Leak} (verbatim survival), \emph{Boundary Error} (partial token survival at word level, length~$>3$), \emph{Type Confusion} (replacement does not match expected format for the entity type), \emph{Format Break} (structured replacement breaks document format), \emph{Over-Masking} (non-entity content unnecessarily altered), \emph{Context Retention} ($>$80\% context word overlap in a $\pm$50\,character window), or \emph{Clean}.
In practice, Type Confusion and Over-Masking produced zero counts across all tested systems, so results tables report the five active categories.
This taxonomy enables targeted diagnosis of where and how systems fail, moving beyond aggregate metrics.

% ============================================================
\section{Experimental Setup}\label{sec:setup}

\subsection{Dataset}

\begin{table}[t]
\centering
\caption{Dataset statistics.}\label{tab:dataset}
\small
\begin{tabular}{@{}lrrrr@{}}
\toprule
Split & Records & Entities & Avg/Rec & PII\% \\
\midrule
Train & 113{,}133 & 286{,}542 & 2.53 & 14.9 \\
Validation & 3{,}600 & 9{,}211 & 2.56 & 15.0 \\
Test & 3{,}600 & 9{,}271 & 2.58 & 15.2 \\
\bottomrule
\end{tabular}
\end{table}

Table~\ref{tab:dataset} summarizes the dataset.
The original texts are drawn from the English subset of the AI4Privacy corpus, which contains crowd-sourced text with synthetic PII.
Entity annotations and anonymized text pairs are produced by the SAHA-AL annotation pipeline: spaCy \texttt{en\_core\_web\_lg} NER and 23~regex patterns pre-annotate entities, a merger resolves overlapping spans with confidence scoring, and Faker generates format-preserving replacements.
Human annotators review and correct entity labels and replacements through a Streamlit interface, producing gold-standard parallel pairs.
The final dataset is augmented via entity swapping ($4\times$), template filling, and Easy Data Augmentation ($3\times$).

\paragraph{Entity type mapping.}
The dataset uses 12~entity type labels, dominated by FULLNAME (5{,}332 entities, 57.5\% of the test set) and UNKNOWN (1{,}813, 19.6\%).
However, the evaluation scripts \emph{re-infer} entity types from text patterns using cascaded heuristics (regex for structured types, rule-based classification for names and organizations), producing a different 13-type mapping.
Under this evaluation mapping, FULLNAME drops to 4{,}177 entities (some reclassified as FIRST\_NAME or absorbed into other categories), while ORGANIZATION grows from~1 to 1{,}062 and a new LOCATION type (753) appears.
The aggregate ELR is unaffected by this re-mapping because it uses text-based matching, but per-type breakdowns (Table~\ref{tab:pertype}) use the evaluation mapping.
Train/test entity string overlap is 43.9\%.

\subsection{Baseline Systems}

\paragraph{Rule-based.}
\emph{Regex+Faker}: ordered regex patterns with Faker-generated replacements.
\emph{spaCy+Faker}~\citep{spacy2020}: \texttt{en\_core\_web\_lg} NER overlaid with regex, Faker replacement.
\emph{Presidio}~\citep{presidio2023}: Microsoft's production SDK with default recognizers and anonymization operators.

\paragraph{Seq2Seq.}
Five encoder-decoder models fine-tuned with \texttt{PIIAwareLoss}, a custom loss function that combines label-smoothed cross-entropy on non-PII tokens with an anti-leakage penalty $-\log(1-p(\text{original token}))$ at PII-span positions, weighted by entity-type sensitivity ($\alpha\!=\!0.3$, $\beta\!=\!2.0$, label smoothing~$0.1$).
All five checkpoints are from this training regime (\texttt{train3.py} in the \texttt{Seq2Seq\_model} codebase; checkpoints stored under \texttt{checkpoints\_pii\_aware\_loss/}).
The models are: BART-base~\citep{lewis2020bart} (139M), Flan-T5-small~\citep{chung2022flant5} (77M), T5-small~\citep{raffel2020t5} (60M), DistilBART (230M), and T5-efficient-tiny (16M).
A sixth model (Flan-T5-base with QLoRA) is implemented but was not included in this evaluation.
Benchmark inference uses beam search ($B\!=\!4$, max length 128).
Note that the \texttt{Seq2Seq\_model} internal evaluation script uses greedy decoding ($B\!=\!1$) on the full test set for speed; all results reported here use $B\!=\!4$.

\subsection{Evaluation Protocol}

All evaluations use the frozen test split (3{,}600 records, 9{,}271 entities).
ERA evaluates 1{,}261 entities---the subset for which typed candidate pools (capped at 500 per type from training data) contain at least two candidates.
LRR evaluates 300 sampled records.
A bootstrap resampling procedure~\citep{efron1993} ($B\!=\!1{,}000$, 95\%) is implemented for key metrics (ELR, token recall, OMR, FPR, CRR-3) but was not executed for the results reported here; we plan to include confidence intervals in a future revision.

\paragraph{Curated test set.}
A separate manually curated test set was constructed via the SAHA-AL annotation pipeline: annotators reviewed and corrected entity labels and replacement quality on a held-out subset of AI4Privacy records, using the Streamlit interface with confidence-based routing into green (high-agreement), yellow, and red (low-confidence) queues.
This set is intended to validate whether results on the synthetic test split generalize to human-verified annotations.
The full benchmark suite has not yet been run on this set.

% ============================================================
\section{Results}\label{sec:results}

\subsection{End-to-End vs.\ Detect-Then-Replace}\label{sec:res-main}

\begin{table}[t]
\centering
\caption{Task~2: Anonymization quality. ELR: entity leakage rate (\%); TokRec: token recall (\%); BERT: BERTScore $F_1$.}\label{tab:anon}
\small
\begin{tabular}{@{}lcccc@{}}
\toprule
System & ELR$\downarrow$ & TokRec$\uparrow$ & OMR$\downarrow$ & BERT$\uparrow$ \\
\midrule
BART-base & \textbf{0.93} & \textbf{95.50} & \textbf{0.13} & 92.74 \\
Flan-T5-sm & 0.99 & 95.33 & 0.22 & 92.47 \\
DistilBART & 1.23 & 94.41 & 0.15 & 86.34 \\
T5-small & 1.54 & 94.15 & 0.19 & 92.59 \\
T5-eff-tiny & 4.14 & 90.42 & 0.49 & 92.57 \\
\midrule
spaCy+Faker & 26.44 & 74.59 & 2.47 & 91.86 \\
Presidio & 33.77 & 68.18 & 1.18 & 90.04 \\
Regex+Faker & 83.39 & 23.05 & 0.00 & 98.15 \\
\bottomrule
\end{tabular}
\end{table}

Table~\ref{tab:anon} reports anonymization quality.
All seq2seq models achieve ELR below 5\%.
BART-base leaks 86 of 9{,}271 entities (0.93\%) while preserving 95.5\% of non-PII tokens at BERTScore $F_1 = 92.74$.
The best rule-based system (spaCy+Faker) leaks 2{,}451 entities (26.44\%)---28$\times$ more.
Regex+Faker leaks 83.39\%.

This is an order-of-magnitude difference, not an incremental improvement.
The gap reflects a qualitative shift: seq2seq models bypass the detection bottleneck by learning implicit entity models during training, while rule-based systems leak every entity the detector misses.

\subsection{Detection is the Bottleneck}\label{sec:res-det}

\begin{table}[t]
\centering
\caption{Task~1: Span-level detection $F_1$.}\label{tab:detection}
\small
\begin{tabular}{@{}lccc@{}}
\toprule
System & Exact & Partial & Type-aware \\
\midrule
spaCy+Faker & 56.30 & 69.38 & 18.18 \\
Presidio & 48.57 & 60.11 & 4.97 \\
Regex+Faker & 25.86 & 27.99 & 24.39 \\
\bottomrule
\end{tabular}
\end{table}

No rule-based detector exceeds 57\% exact $F_1$ (Table~\ref{tab:detection}).
Structured types (EMAIL, CREDIT\_CARD, SSN) are detected at $>$92\% recall by all systems.
FULLNAME recall ranges from 0\% (Regex) to 39\% (spaCy); ADDRESS recall is 0\% across all systems.

In a detect-then-replace pipeline, every missed entity is leaked verbatim.
spaCy misses 4{,}863 of 9{,}271 entities; of its 2{,}451 leaked entities, the overwhelming majority are undetected names.
Detection recall is the ceiling on rule-based anonymization quality, and that ceiling is low.

\subsection{Privacy Under Attack}\label{sec:res-attack}

\begin{table}[t]
\centering
\caption{Task~3: Adversarial privacy risk.}\label{tab:privacy}
\small
\begin{tabular}{@{}lcccc@{}}
\toprule
System & CRR-3$\downarrow$ & ERA@1$\downarrow$ & ERA@5$\downarrow$ & UAC$\downarrow$ \\
\midrule
BART-base & \textbf{34.62} & \textbf{1.90} & \textbf{4.84} & \textbf{0.33} \\
Flan-T5-sm & 34.94 & 1.67 & 5.08 & 0.22 \\
spaCy+Faker & 40.62 & 9.36 & 14.27 & 1.58 \\
Presidio & 50.33 & 20.46 & 26.09 & 1.78 \\
\bottomrule
\end{tabular}
\end{table}

\begin{table}[t]
\centering
\caption{Attack comparison for BART-base.}\label{tab:attack_compare}
\small
\begin{tabular}{@{}lrl@{}}
\toprule
Attack & Recovery & Method \\
\midrule
LRR exact (Qwen 14B) & 0.13\% & Generative inference \\
LRR fuzzy (Qwen 14B) & 0.53\% & Generative (fuzzy) \\
ELR & 0.93\% & Verbatim string match \\
ERA@1 & 1.90\% & Embedding retrieval, top-1 \\
ERA@5 & 4.84\% & Embedding retrieval, top-5 \\
\bottomrule
\end{tabular}
\end{table}

Table~\ref{tab:privacy} reports adversarial metrics.
BART-base leads on all measures.
Table~\ref{tab:attack_compare} reveals the attack effectiveness hierarchy:
\[
\text{LRR}_{\text{exact}} (0.13\%) < \text{ELR} (0.93\%) < \text{ERA@1} (1.90\%) < \text{ERA@5} (4.84\%).
\]

Retrieval-based attacks recover $\approx$15$\times$ more entities than generative LLM attacks (1.90\% vs.\ 0.13\%).
This is the central finding of the adversarial evaluation: a system that appears to achieve $<$1\% leakage under passive observation is exposed at nearly 5\% under a retrieval adversary with modest auxiliary knowledge.
Standard ELR-only evaluation would miss this entirely.

\subsection{Where Models Fail}\label{sec:res-fail}

\begin{table}[t]
\centering
\caption{Per-type ELR for BART-base.}\label{tab:pertype}
\small
\begin{tabular}{@{}lrrl@{}}
\toprule
Type & Leaked & Total & ELR \\
\midrule
FULLNAME & 45 & 4{,}177 & 1.08\% \\
UNKNOWN & 26 & 1{,}081 & 2.41\% \\
LOCATION & 6 & 753 & 0.80\% \\
ORGANIZATION & 6 & 1{,}062 & 0.57\% \\
TIME & 2 & 337 & 0.59\% \\
DATE & 1 & 599 & 0.17\% \\
EMAIL/PHONE/SSN/CC/ID/ZIP & 0 & 1{,}262 & 0.00\% \\
\bottomrule
\end{tabular}
\end{table}

\begin{table}[t]
\centering
\caption{Failure taxonomy (\% of 9{,}271 entities).}\label{tab:taxonomy}
\small
\begin{tabular}{@{}lcccc@{}}
\toprule
Category & BART & spaCy & Presid. & Regex \\
\midrule
Clean & 61.5 & 42.0 & 28.2 & 3.5 \\
Full Leak & 0.9 & 26.4 & 33.8 & 83.4 \\
Ctx Retention & 36.4 & 29.4 & 25.2 & 12.7 \\
Boundary Err & 1.1 & 1.8 & 2.1 & 0.3 \\
Format Break & 0.1 & 0.4 & 10.7 & 0.0 \\
\bottomrule
\end{tabular}
\end{table}

Structured PII is effectively eliminated under the tested (synthetic) distribution.
All seq2seq models achieve 0\% leakage on EMAIL, PHONE, SSN, CREDIT\_CARD, ID\_NUMBER, and ZIPCODE (Table~\ref{tab:pertype}); whether this holds on real-world data with diverse international formats remains an open question.
Names remain the open problem: FULLNAME and UNKNOWN account for 71 of BART-base's 86 leaks (82.6\%).

The failure taxonomy (Table~\ref{tab:taxonomy}) reveals qualitatively different error profiles.
BART-base's dominant non-clean category is context retention (36.4\%), which reflects faithful non-PII preservation---not re-identification risk.
Its actual privacy failures (Full Leak + Boundary Error) total 2.0\%.
Presidio's distinctive failure is format breaks (10.7\%): its default operator emits tags like \texttt{<DATE\_TIME>} instead of natural replacements, destroying downstream utility.
Rule-based systems exhibit systematic, categorical failure (83.4\% full leak for Regex); seq2seq systems produce rare, specific failures concentrated on unusual names.

\subsection{Pareto Frontier}\label{sec:res-pareto}

\begin{table}[t]
\centering
\caption{Privacy-utility tradeoff. Privacy $= 1 - \text{ELR}/100$; Utility $= \text{BERTScore}/100$.}\label{tab:pareto}
\small
\begin{tabular}{@{}lccc@{}}
\toprule
System & Privacy & Utility & PUS$_{0.5}$ \\
\midrule
BART-base & \textbf{0.991} & 0.927 & \textbf{0.959} \\
Flan-T5-sm & 0.990 & 0.925 & 0.957 \\
T5-small & 0.985 & 0.926 & 0.955 \\
T5-eff-tiny & 0.959 & 0.926 & 0.942 \\
DistilBART & 0.988 & 0.863 & 0.926 \\
spaCy+Faker & 0.736 & 0.919 & 0.827 \\
Presidio & 0.662 & 0.900 & 0.781 \\
Regex+Faker & 0.166 & 0.982 & 0.574 \\
\bottomrule
\end{tabular}
\end{table}

Two systems lie on the Pareto frontier (Table~\ref{tab:pareto}): Regex+Faker (highest utility, lowest privacy) and BART-base (highest privacy with competitive utility).
All other systems are strictly dominated on both axes, though several seq2seq variants (Flan-T5-small, T5-small) fall within a narrow margin of BART-base.
BART-base is the only system whose privacy score (0.991) exceeds its utility score (0.927), so its PUS increases with $\lambda$ and it ranks first for any $\lambda \geq 0.07$.

% ============================================================
\section{Discussion}\label{sec:discussion}

\paragraph{Why end-to-end models succeed.}
Detect-then-replace pipelines expose a hard coupling: anonymization quality is bounded by detection recall, and undetected entities leak with probability one.
Seq2seq models decouple detection from replacement by learning to rewrite text holistically.
During training, the model implicitly learns which tokens constitute PII and how to generate contextually consistent replacements, without requiring an explicit entity inventory.
This explains why BART-base achieves 0.93\% ELR despite never being trained as an NER system: the entity model is latent in the encoder-decoder weights.

\paragraph{Why BERTScore alone is misleading.}
Regex+Faker achieves the highest BERTScore (98.15)---seemingly the best utility---while leaking 83.39\% of entities.
The explanation is trivial: Regex detects only 15\% of entities, so 85\% of the text is unchanged, yielding near-perfect semantic similarity.
A system that modifies nothing achieves BERTScore~$\approx$\,100 and ELR~$\approx$\,100.
BERTScore~\citep{zhang2020bertscore} measures utility but is blind to privacy.
The PUS score correctly penalizes this failure: Regex ranks last at PUS\,=\,0.574.
This underscores the need for composite evaluation---no single metric captures both privacy and utility.

\paragraph{Why retrieval attacks matter more than generative ones.}
Replacement-based anonymization substitutes one plausible name for another, producing contextually coherent text.
A generative adversary (LLM) cannot distinguish ``Michael Jones'' from ``John Smith'' based on surrounding context, because both are valid completions---hence LRR~$< 0.2\%$.
A retrieval adversary does not reason about plausibility.
It exploits statistical correlations between the anonymized text's global embedding and candidate entity embeddings, bypassing the semantic consistency of the replacement.
ERA@1 recovers 1.90\% of entities precisely because embedding-level signal persists after surface-level replacement.

The implication is direct: \textbf{surface anonymization does not equal semantic privacy.}
Text that passes all standard quality metrics may still be vulnerable to an adversary with embedding access and auxiliary knowledge.
Future anonymization systems must consider embedding-level defenses, not just surface-level text replacement.

\paragraph{Model capacity and architecture.}
ELR scales approximately log-linearly with parameter count from T5-eff-tiny (16M, 4.14\%) to BART-base (139M, 0.93\%).
DistilBART (230M) breaks this trend with 1.23\% ELR but only 86.34 BERTScore, likely because its summarization-oriented pretraining produces more aggressive text rewriting.
Architecture and pretraining objective matter as much as raw parameter count.

% ============================================================
\section{Limitations}\label{sec:limitations}

This section warrants careful reading.

\paragraph{Pending evaluation on curated test set.}
The manually curated test set (Section~\ref{sec:setup}) has been constructed but the full benchmark suite has not yet been run on it.
Reported ELR results may be optimistic relative to human-verified annotations.
A rigorous ablation on this dataset is required to validate robustness claims.

\paragraph{Synthetic data.}
Original texts are drawn from the AI4Privacy corpus, which uses crowd-sourced templates with Faker-generated PII.
Entity distributions, naming patterns, and contextual structure may not reflect real-world PII distributions (medical records, legal documents, social media).
Results should be interpreted as a controlled comparison, not as deployment-ready estimates.

\paragraph{Train/test entity overlap.}
43.9\% of unique test entity strings appear in training data, a consequence of Faker's finite name/email pools.
Seq2seq models may partially benefit from memorization.
Absolute ELR figures are potentially inflated; a zero-overlap evaluation split would be needed to disentangle memorization from generalization.

\paragraph{Limited adversary models.}
ERA uses a closed-world candidate pool (training-set entities, 500 per type).
Real adversaries may possess broader auxiliary knowledge.
LRR uses a fixed prompt (Appendix~\ref{app:lrr}); adversarial prompt engineering or white-box access to the anonymization model could yield higher recovery rates.
The current threat model represents a lower bound on adversarial capability.

\paragraph{Entity type ambiguity.}
Approximately 19.6\% of test entities are typed as UNKNOWN due to heuristic inference limitations, reducing the reliability of per-type analysis.

\paragraph{English only.}
All data, models, and evaluation are English-specific.

\paragraph{No formal privacy guarantees.}
The benchmark provides empirical privacy metrics.
ELR~$<$\,1\% does not constitute a differential privacy guarantee~\citep{dwork2006} or imply regulatory compliance.

\paragraph{No efficiency evaluation.}
Inference latency, memory, and computational cost are not measured.

% ============================================================
\section{Future Work}\label{sec:future}

\begin{enumerate}
\item \textbf{Curated test set evaluation.}
Run the full benchmark and ablation study on the manually curated test set to validate ELR robustness and quantify performance degradation relative to the synthetic test split.
This is the most critical next step.
\item \textbf{Stronger adversaries.}
Evaluate adaptive adversaries with white-box access to the anonymization model, adversarial prompt engineering for LRR, and open-world candidate pools for ERA.
\item \textbf{Embedding-level privacy defenses.}
Investigate noise injection in the encoder latent space, structurally diverse replacement generation (different name length, ethnicity, gender), and representation scrubbing to reduce ERA vulnerability.
\item \textbf{Real-world datasets.}
Extend evaluation to clinical notes, legal documents, and social media text to assess generalization beyond synthetic templates.
\item \textbf{Zero-overlap split.}
Construct an evaluation split with zero train/test entity string overlap to measure generalization without memorization confounds.
\item \textbf{Formal privacy integration.}
Combine empirical SAHA-AL metrics with differentially private text generation mechanisms to bridge the gap between empirical and theoretical guarantees.
\end{enumerate}

% ============================================================
\section{Conclusion}\label{sec:conclusion}

SAHA-AL provides an evaluation framework that treats text anonymization as a system under adversarial attack rather than a standalone text transformation.
On a synthetic test set, fine-tuned seq2seq models reduce entity leakage from 26--83\% (rule-based) to below 1\%, an order-of-magnitude improvement.
However, retrieval-based adversaries recover 1.9\% of entities even from the best-performing system---a residual risk invisible to standard leakage metrics.
Structured PII leakage is effectively eliminated under the tested synthetic distribution; names remain the dominant source of failure.

These findings are subject to important caveats: the data is synthetic, train/test overlap is substantial, and evaluation on a separately curated test set is pending.
The central claim stands regardless: anonymization benchmarks that measure only detection accuracy or surface-level leakage are insufficient.
Privacy evaluation must incorporate adversarial threat models.

% ============================================================
\appendix
\section{LRR Prompt}\label{app:lrr}

The following prompt is used for the LLM Re-identification Rate (LRR) evaluation. It is passed as a single user message to the generative adversary (Qwen~2.5-14B or Llama~3.3-70B). The \texttt{\{anonymized\_text\}} placeholder is replaced with the anonymized output for each record.

\begin{quote}
\small\ttfamily
You are given an anonymized text where personal information has been replaced with fake alternatives.

Anonymized text: \{anonymized\_text\}

For each piece of replaced information, try to guess the ORIGINAL value based on contextual clues. Output one guess per line in the format:\\
REPLACEMENT -> YOUR\_GUESS

Only output guesses, nothing else.
\end{quote}

Responses are parsed line-by-line; each line containing ``\texttt{->}'' yields a guess from the right-hand side. An entity is counted as an exact match if it appears verbatim among the parsed guesses, and as a fuzzy match if any guess achieves $>0.8$ character-level similarity (Python \texttt{SequenceMatcher} ratio). For the API path, \texttt{temperature=0} and \texttt{max\_tokens=512}; for local models, \texttt{temperature=0.1}, \texttt{top\_p=0.9}, \texttt{max\_new\_tokens=256}.

% ============================================================
\bibliographystyle{plainnat}
\begin{thebibliography}{16}

\bibitem[Lison et~al.(2021)]{lison2021}
Lison, P., Pilán, I., Sanchez, D., Batet, M., \& {\O}vrelid, L. (2021).
\newblock Anonymisation models for text data: State of the art, challenges and future directions.
\newblock In \emph{Proceedings of the 59th Annual Meeting of the Association for Computational Linguistics (ACL)}, pp.\ 4188--4203.

\bibitem[Narayanan \& Shmatikov(2008)]{narayanan2008}
Narayanan, A. \& Shmatikov, V. (2008).
\newblock Robust de-anonymization of large sparse datasets.
\newblock In \emph{2008 IEEE Symposium on Security and Privacy (S\&P)}, pp.\ 111--125.

\bibitem[Pilán et~al.(2022)]{pilan2022tab}
Pilán, I., Lison, P., {\O}vrelid, L., Papadopoulou, A., Sánchez, D., \& Batet, M. (2022).
\newblock The Text Anonymization Benchmark ({TAB}): A dedicated corpus and evaluation framework for text anonymization.
\newblock \emph{Computational Linguistics}, 48(4):1053--1101.

\bibitem[Stubbs et~al.(2015)]{stubbs2015}
Stubbs, A., Kotfila, C., \& Uzuner, {\"O}. (2015).
\newblock Automated systems for the de-identification of longitudinal clinical narratives: Overview of 2014 i2b2/{UT}Health shared task Track~1.
\newblock \emph{Journal of Biomedical Informatics}, 58:S11--S19.

\bibitem[Chung et~al.(2022)]{chung2022flant5}
Chung, H.~W., Hou, L., Longpre, S., Zoph, B., Tay, Y., Fedus, W., Li, Y., Wang, X., Dehghani, M., Brahma, S., et~al. (2022).
\newblock Scaling instruction-finetuned language models.
\newblock \emph{arXiv preprint arXiv:2210.11416}.

\bibitem[Dwork(2006)]{dwork2006}
Dwork, C. (2006).
\newblock Differential privacy.
\newblock In \emph{Proceedings of the 33rd International Colloquium on Automata, Languages and Programming (ICALP)}, pp.\ 1--12. Springer.

\bibitem[Efron \& Tibshirani(1993)]{efron1993}
Efron, B. \& Tibshirani, R.~J. (1993).
\newblock \emph{An Introduction to the Bootstrap}.
\newblock Chapman \& Hall/CRC.

\bibitem[Honnibal et~al.(2020)]{spacy2020}
Honnibal, M., Montani, I., Van~Landeghem, S., \& Boyd, A. (2020).
\newblock {spaCy}: Industrial-strength natural language processing in {P}ython.
\newblock \url{https://spacy.io}.

\bibitem[Lewis et~al.(2020)]{lewis2020bart}
Lewis, M., Liu, Y., Goyal, N., Ghazvininejad, M., Mohamed, A., Levy, O., Stoyanov, V., \& Zettlemoyer, L. (2020).
\newblock {BART}: Denoising sequence-to-sequence pre-training for natural language generation, translation, and comprehension.
\newblock In \emph{Proceedings of the 58th Annual Meeting of the Association for Computational Linguistics (ACL)}, pp.\ 7871--7880.

\bibitem[Microsoft(2023)]{presidio2023}
Microsoft (2023).
\newblock Presidio: Data protection and de-identification {SDK}.
\newblock \url{https://github.com/microsoft/presidio}.

\bibitem[Raffel et~al.(2020)]{raffel2020t5}
Raffel, C., Shazeer, N., Roberts, A., Lee, K., Narang, S., Matena, M., Zhou, Y., Li, W., \& Liu, P.~J. (2020).
\newblock Exploring the limits of transfer learning with a unified text-to-text transformer.
\newblock \emph{Journal of Machine Learning Research}, 21(140):1--67.

\bibitem[Reimers \& Gurevych(2019)]{reimers2019}
Reimers, N. \& Gurevych, I. (2019).
\newblock {Sentence-BERT}: Sentence embeddings using {S}iamese {BERT}-networks.
\newblock In \emph{Proceedings of the 2019 Conference on Empirical Methods in Natural Language Processing (EMNLP)}, pp.\ 3982--3992.

\bibitem[Shokri et~al.(2017)]{shokri2017}
Shokri, R., Stronati, M., Song, C., \& Shmatikov, V. (2017).
\newblock Membership inference attacks against machine learning models.
\newblock In \emph{2017 IEEE Symposium on Security and Privacy (S\&P)}, pp.\ 3--18.

\bibitem[Staab et~al.(2024)]{staab2024}
Staab, R., Vero, M., Balunovi{\'c}, M., \& Vechev, M. (2024).
\newblock Beyond memorization: Violating privacy via inference with large language models.
\newblock In \emph{International Conference on Learning Representations (ICLR)}.

\bibitem[Sweeney(2002)]{sweeney2002}
Sweeney, L. (2002).
\newblock $k$-Anonymity: A model for protecting privacy.
\newblock \emph{International Journal of Uncertainty, Fuzziness and Knowledge-Based Systems}, 10(5):557--570.

\bibitem[Zhang et~al.(2020)]{zhang2020bertscore}
Zhang, T., Kishore, V., Wu, F., Weinberger, K.~Q., \& Artzi, Y. (2020).
\newblock {BERT}Score: Evaluating text generation with {BERT}.
\newblock In \emph{International Conference on Learning Representations (ICLR)}.

\end{thebibliography}

\end{document}